\section{算法设计}\label{sec:algo}

整个求解是确定性的（无随机数），由六个阶段组成，全部针对向量化的候选点数组批量计算。

\subsection{不确定集与误差网格离散}

源不确定集按 $d$（$31$ 点，步长 $50$ m，含端点 $5$ m 与 $1500$ m）$\times$ $\delta_1$
（$9$ 点，步长 $0.25^\circ$，含 $\pm1^\circ$）离散为 $279$ 个采样点，四个极点必在其中；
第二次测向误差 $\delta_2$ 取 $5$ 点（步长 $0.5^\circ$，含 $\pm1^\circ$ 两端）。
于是每个候选点 $S_2$ 的代价函数是 $279\times5=1395$ 个楔形交直径的最大值，全部用数组
运算一次算完。端点必须在网格内——最坏情形恰好出现在 $d=1500$ m 与 $\delta_1=\delta_2=+1^\circ$
的端点上。

\subsection{全域粗搜与局部细化}

在 $1800$ m 圆域上取 $25$ m 直角网格（$21\,025$ 格）。由于 $J$ 只对可行点有意义，先算每格
到源采样集的最坏距离、按式~\eqref{eq:feasible} 筛出可行格（$712$ 个，占 $3.4\%$），
\textbf{再只对可行格计算 $J$}，省去 $96\%$ 以上的无效计算，取其中最小者作为粗解。
随后以粗解为中心、$\pm80$ m 范围、$2$ m 步长再搜一遍，把最优值定到米级（$2$ m 的偏差远小于
$\pm1^\circ$ 测向误差对应的几十米量级）。

\subsection{文献快筛与全局认证}

粗搜解依赖"粗网格是否恰好框住最优盆地"。为避免这个风险，利用 §\ref{sec:lit} 的结论
（GDOP 与精确 $J$ 的秩相关 $0.998$）设计一条独立路线：

\begin{enumerate}
  \item 在可行域的包围盒内取 $5$ m 细网格，筛出全部可行点（$17\,861$ 个）；
  \item 用解析 GDOP 场（式~\eqref{eq:gdop}，含全体源采样点取最坏）对这些点打分，
        取分数最优的 $48$ 个候选（含当前粗搜解）；
  \item 对候选做精确构造复核，再对最优候选做 $2$ m 局部细化；细化结果与粗搜细化结果取优者。
\end{enumerate}

实测两条路线给出 $J=134.0831$ m，彼此相距 $8.49$ m、最坏直径相差 $0.063$ m
（细网格路线更优），说明粗网格没有漏掉更好的盆地；本文最终采用细网格认证解。
当可行域明显变大时（例如更换第一检测点后），程序按点数上限自动放宽网格步长并如实记录，
保证换参数时不会出现网格爆炸——三组不同检测点的实测运行时间为 $15\sim17$ s。

\subsection{候选弧带的构造}

候选区域按式~\eqref{eq:cand} 在极坐标下构造：绕 $S_1$ 取径向步长 $5$ m、方位步长
$0.5^\circ$ 的网格，对每个方位角求满足 $J\le(1+\eta)J^*$ 且可行的\textbf{连续}径向区间
（实测每个方位角上区间均连续、无断口），再把逐方位的径向下界／上界组装成弧带多边形。
边界口径为：若紧邻的下一个网格点已不可行，则该侧直接取可行域透镜的\textbf{精确射线距离}
（解析式；最优解恰好落在这条边界上，故必须用精确值收边），否则取网格单元中心外扩／内缩
半个径向步长。因此报告的范围与面积与真实次水平集的差别不超过半个网格单元
（径向 $2.5$ m、方位 $0.25^\circ$）。

\subsection{镜像规范与并列解}

源不确定集、误差网格、圆域与代价函数在关于示向度方向的\textbf{反射}下都不变，因此若
$S_2$ 最优，则其镜像点也最优，两瓣候选区域面积严格相等（实测比 $1.0000$）。程序把
"$\varphi<0$ 的解反射到 $\varphi>0$"作为报告规范，并每次运行核对镜像点代价与最优值一致
（实测相对偏差 $1\times10^{-14}$ 以内）。

\subsection{计算代价}

全流程约 $15$ s（含两次抽样校验），主要开销是 $1395$ 组合的最坏直径向量化计算；解析 GDOP
场的引入使细网格快筛的开销可以忽略。六个产物（两图、两矢量图、逐格适合度表、结论与校验
JSON）在两次运行之间逐字节一致，即整套结果可完全复现。
